\chapter{A Minimal Accessible Preamble}
\label{app:minimal}

You do not need this whole class to make an accessible document. The
skeleton below is the smallest starting point: the metadata line
(PDF/UA-2, so formulas enter the tag tree as readable MathML), the Lua
font setup with a sans-serif body face, and the two content rules
(alt text, table headers). Compile it with LuaLaTeX.

% commandchars typesets the leading backslashes via |textbackslash() so the
% source has no literal class/metadata command; the printed example is
% identical, but Overleaf will not mistake this appendix for a "main document".
\begin{Verbatim}[numbers=left,commandchars=\|\(\)]
|textbackslash()DocumentMetadata{lang=en-US,
  pdfstandard=ua-2, tagging=on,
  tagging-setup={math/setup={mathml-AF,mathml-SE}}}
|textbackslash()documentclass{article}
\usepackage{fontspec}      % LuaLaTeX fonts
\setmainfont{Arial}        % sans-serif body font
\usepackage{unicode-math}  % OpenType math
\usepackage{graphicx}
\usepackage{booktabs}
\begin{document}

\section{A heading}
Some text, and math $e^{i\pi}+1=0$.

\begin{figure}\centering
  \includegraphics[width=.5\linewidth,
    alt={Describe the image}]{pic}
  \caption{Why this figure is here.}
\end{figure}

\tagpdfsetup{table/header-rows={1}}
\begin{tabular}{lr}
  \toprule
  Name & Value \\
  \midrule
  A & 1 \\
  \bottomrule
\end{tabular}

\end{document}
\end{Verbatim}


\section{Where to read more}
\label{sec:more}

The accessibility machinery used here comes from the \LaTeX\ Tagging
Project; for the authoritative, continually updated instructions, see
\hrefu{https://latex3.github.io/tagging-project/}{the Tagging Project
documentation}. For another university template built on the same
machinery -- PDF/UA-2, tagged tables, and screen-readable mathematics
-- see \hrefu{https://github.com/graduatecollege/uofithesis}{the
University of Illinois Graduate College's \texttt{uofithesis} class}.
And for a broader collection of \LaTeX\ tutorials, reference sheets,
and accessibility-minded thesis, homework, and poster templates, see
\hrefu{https://pauljhurtado.com/latex}{Paul Hurtado's \LaTeX\ resource
page}.